% Generated draft: requires mathematical and source review.
\documentclass[11pt]{article}
\usepackage[margin=25mm]{geometry}
\usepackage{fontspec}
\setmainfont{Latin Modern Roman}
\usepackage{amsmath,amssymb,mathtools,unicode-math}
\setmathfont{Latin Modern Math}
\usepackage{enumitem,needspace,xurl,hyperref}
\setlength{\parindent}{0pt}
\setlength{\parskip}{5pt}
\setlength{\emergencystretch}{3em}
\newcommand{\N}{\mathbb{N}}
\newcommand{\Z}{\mathbb{Z}}
\newcommand{\Q}{\mathbb{Q}}
\newcommand{\R}{\mathbb{R}}
\newcommand{\C}{\mathbb{C}}
\newcommand{\F}{\mathbb{F}}
\newcommand{\A}{\mathbb{A}}
\newcommand{\PP}{\mathbb{P}}
\newcommand{\E}{\mathbb{E}}
\newcommand{\eps}{\varepsilon}
\newcommand{\dd}{\,\mathrm{d}}
\begin{document}
\section*{UnsolvedMath ID: 30003447}
\textbf{Source problem code:} catalogue entry
\subsection{Definitions and mathematical statement}
The supplied record leaves the domain of positivity and the precise meaning of flatness implicit. The following gives a conventional real, unconstrained interpretation; these conventions require confirmation against the source.

Fix integers \(n\geq 1\) and \(m\geq 0\). For \(\alpha\in\Z_{\geq 0}^{n}\), put \(|\alpha|=\alpha_1+\cdots+\alpha_n\) and \(x^\alpha=x_1^{\alpha_1}\cdots x_n^{\alpha_n}\). Write \(\R[x]_{\leq m}\) for the real polynomials in \(n\) variables of total degree at most \(m\). A real truncated moment sequence is an array
\[
\beta^{(m)}=(\beta_\alpha)_{|\alpha|\leq m},\qquad \beta_\alpha\in\R.
\]
Its Riesz functional is the linear map
\[
L_\beta:\R[x]_{\leq m}\longrightarrow\R,\qquad
L_\beta\!\left(\sum_{|\alpha|\leq m}p_\alpha x^\alpha\right)
=\sum_{|\alpha|\leq m}p_\alpha\beta_\alpha.
\]
Under the full-space interpretation, positivity means
\[
L_\beta(p)\geq 0
\quad\text{for every }p\in\R[x]_{\leq m}
\text{ such that }p(x)\geq 0\text{ for every }x\in\R^n.
\]
This condition must be distinguished from requiring only \(L_\beta(q^2)\geq 0\) for every \(q\) with \(2\deg q\leq m\). The record does not explicitly specify which meaning is intended.

For an integer \(d\geq 1\) and a sequence \(\beta^{(2d)}\), define the moment matrix
\[
M_d(\beta)=\bigl(\beta_{\alpha+\gamma}\bigr)_{|\alpha|,|\gamma|\leq d}.
\]
Here positive means positive semidefinite, denoted \(M_d(\beta)\succeq 0\). Under the usual adjacent-order rank convention, this matrix is flat when
\[
\operatorname{rank}M_d(\beta)=\operatorname{rank}M_{d-1}(\beta),
\]
where \(M_{d-1}(\beta)\) is its principal submatrix indexed by monomials of degree at most \(d-1\). With this convention, put
\[
\mathcal S_{n,d}=\left\{M_d(\beta):
\beta\in\R^{\binom{n+2d}{n}},\ M_d(\beta)\succeq 0,
\ \operatorname{rank}M_d(\beta)=\operatorname{rank}M_{d-1}(\beta)\right\},
\qquad
F_{n,d}=\overline{\mathcal S_{n,d}},
\]
where the closure is taken in the finite-dimensional real vector space of moment matrices, with its Euclidean topology. Equivalently, \(A\in F_{n,d}\) means that there exists a sequence \(A_k\in\mathcal S_{n,d}\) converging entrywise to \(A\).

\textbf{Problem, subject to the stated conventions.} Give concrete conditions on \(\beta^{(m)}\) sufficient to ensure positivity of \(L_\beta\), and give a concrete necessary-and-sufficient characterization of membership in \(F_{n,d}\) for general \(n\) and \(d\).

The source calls the flat positive matrices a cone. The definition above uses the set of flat positive moment matrices itself; whether the source instead intends a conic hull must be checked. The source also does not specify normalization, a support set, the relation of \(m\) to \(2d\) beyond the matrix construction, or boundary conventions for \(d=0\). No such additional hypotheses are imposed here.
\subsection{Short English statement}
Find concrete sufficient conditions ensuring that a truncated moment sequence has a positive Riesz functional, and characterize limits of flat positive moment matrices. The displayed formulation uses explicit conventional definitions that need source confirmation. The supplied status is partially solved; this is a statement-only draft requiring human mathematical review.
\subsection{Sources}
\begin{itemize}
\item Supplied local UnsolvedMath registry record, permanent ID \(30003447\), exact source problem code \texttt{catalogue entry}, title: Positivity and Closure of Flat Moment Matrices. Collection: Oberwolfach Reports Open Problems. The record gives proposed year 2017 and names no individual proposer.
\item Supplied source citation: \emph{Real Algebraic Geometry With a View Toward Moment Problems and Optimization} (2018), Oberwolfach Reports, pp.\ 771--862, DOI 10.4171/owr/2017/14. Original Oberwolfach report: \url{https://doi.org/10.4171/owr/2017/14}. The supplied background describes these reports as documenting mathematical workshops and their open-problem discussions.
\item Supplied literature reference: \emph{Positive Gorenstein ideals}, \url{https://doi.org/10.1090/S0002-9939-2014-12253-2}. Authors and a precise supporting passage are not supplied in this record.
\item Literature assessment date: \texttt{2026-08-21}. This is supplied dated curation, not independently verified. The supplied assessment labels the problem partially solved and states that concrete characterizations are known in several low-dimensional cases and rank-based positivity criteria are known, while a general concrete characterization of \(F\) remains open. The background describes this assessment as dated literature triage derived from the supplied OWR curation and calls for independent verification of time-sensitive status and cited claims.
\item The registry labels the statement \texttt{corrected\_verified}, explaining that mathematical content was preserved while notation, typography, and sentence structure were repaired for publication. This is a supplied editorial label, not independent mathematical verification of this draft or its conventions.
\end{itemize}
\end{document}
